% !TEX root = ../main.tex

\chapter{Geometric constructions in \texorpdfstring{$\Cinfty$}{C-infinity}-rings}
\label[chapter]{chap:Manifolds_As_Cinfty_Rings}

The preceding chapter showed that a smooth manifold and all of its smooth maps
are encoded by its ring of smooth functions. We now ask whether geometric
constructions are encoded as well. Since the functor
\[
  \Cinfty(-)\colon
  (\text{smooth manifolds})\catop
  \longrightarrow
  \Cinfty\text{-}\mathrm{Ring}
\]
is contravariant, the first test is whether a pullback of manifolds becomes a
pushout of $\Cinfty$-rings.

\begin{question}
  \label[question]{question:When_Pushout_Is_Manifold}
  When is the algebraic pushout of smooth function rings again the function
  ring of the geometric pullback?
\end{question}

Transversality is a sufficient condition. To understand why, we first show
that independent smooth equations generate the full ideal of functions
vanishing on their common zero locus. This lets us translate geometric
constructions into presentations of smooth rings.

Two intermediate results prepare the general pullback theorem. Every
manifold has a finitely presented function ring, and products and open
restrictions have the expected algebraic descriptions. We then prove the
transverse-pullback theorem locally and explain its consequences. The final
proof section supplies the globalization argument. Together with the full
faithfulness proved in the preceding chapter, these results identify
manifolds inside smooth algebra and explain which intersections that
embedding preserves.

\section{Equations and zero loci}
\label[section]{sec:Equations_And_Coproducts}

Return to the intersections in \Cref{fig:Two_Intersections}. Both the
transverse axes and the tangent line and parabola intersect in a single
point. Their equations nevertheless give different algebraic pushouts.
We begin by making this calculation precise, then ask when a quotient by
equations is the function ring of their zero locus.

\subsection{Coproducts impose simultaneous equations}
\label[subsection]{sec:Coproducts_Impose_Equations}

We first fix the algebraic notation. The category of $\Cinfty$-rings has all
small colimits. Following Moerdijk and Reyes, we write
\[
  B\otimes_\infty C
\]
for the coproduct of two $\Cinfty$-rings. This notation records the
categorical role of the construction. It is not, in general, the tensor
product of the underlying commutative $\R$-algebras.

\begin{proposition}[Coproducts of free rings and quotients]
  \label[proposition]{prop:Cinfty_Coproduct_Formulas}
  The following formulas hold.
  \begin{enumerate}
    \item There is a natural isomorphism
      \[
        \Cinfty(\R^n)\otimes_\infty\Cinfty(\R^m)
        \cong \Cinfty(\R^{n+m}).
      \]
    \item If $I\subseteq B$ and $J\subseteq C$ are ideals, then
      \[
        (B/I)\otimes_\infty(C/J)
        \cong
        (B\otimes_\infty C)/(I,J),
      \]
      where $(I,J)$ denotes the ideal generated by the images of $I$ and $J$.
    \item If $I,J\subseteq B$, then
      \[
        (B/I)\mathbin{\amalg_B}(B/J)
        \cong B/(I+J).
      \]
  \end{enumerate}
\end{proposition}

\begin{proof}
  A homomorphism from $\Cinfty(\R^n)\otimes_\infty\Cinfty(\R^m)$ to a
  $\Cinfty$-ring $D$ is a pair of homomorphisms from the two factors. By
  \Cref{thm:Free_Cinfty_Rings}, this is the same as a pair consisting of an
  $n$-tuple and an $m$-tuple of elements of $D$, hence the same as an
  $(n+m)$-tuple. The free-ring theorem identifies the representing object as
  $\Cinfty(\R^{n+m})$.

  A homomorphism from the coproduct in the second formula to $D$ is a pair of
  homomorphisms $B\to D$ and $C\to D$ which kill $I$ and $J$, respectively.
  By the universal properties of the coproduct and the quotient, this is the
  same as a homomorphism from $(B\otimes_\infty C)/(I,J)$ to $D$. The third
  formula is proved in the same way: a compatible pair of maps from $B/I$ and
  $B/J$ is a map from $B$ killing both ideals.
\end{proof}

Set
\[
  A=\Cinfty(\R^2).
\]
For the two coordinate axes, the function rings are $A/(x)$ and $A/(y)$.
Imposing both equations gives
\[
  A/(x)\mathbin{\amalg_A}A/(y)
  \cong A/(x,y)
  \cong \R.
\]
This is the function ring of their intersection point. For the $x$-axis and
the parabola $y=x^2$, the analogous calculation begins
\[
  A/(y)\mathbin{\amalg_A}A/(y-x^2)
  \cong A/(y,y-x^2).
\]
Eliminating $y$ gives $\Cinfty(\R)/(x^2)$, the dual-number ring.

The first calculation recovers the function ring of the geometric
intersection; the second retains a nonzero nilpotent element. The
regular-equation theorem below explains why independent equations
recover the function ring of the zero locus.

\subsection{Regular equations}
\label[subsection]{sec:Regular_Equations}

Let $M$ be a smooth manifold and let $g_1,\ldots,g_q\in\Cinfty(M)$. Write
\[
  Z(g_1,\ldots,g_q)
  =\{p\in M\mid g_1(p)=\cdots=g_q(p)=0\}
\]
and
\[
  \mathfrak m_Z
  =\{h\in\Cinfty(M)\mid h|_Z=0\}.
\]
The inclusion $(g_1,\ldots,g_q)\subseteq\mathfrak m_Z$ is automatic. Equality
is a geometric assertion about the equations.

We first isolate the globalization argument which will occur several times.

\begin{lemma}[Local ideal membership]
  \label[lemma]{lem:Local_Ideal_Membership}
  Let $h,g_1,\ldots,g_q\in\Cinfty(M)$. If every point of $M$ has an open
  neighborhood $U$ such that
  \[
    h|_U\in(g_1|_U,\ldots,g_q|_U),
  \]
  then $h\in(g_1,\ldots,g_q)$ in $\Cinfty(M)$.
\end{lemma}

\begin{proof}
  Choose a locally finite open cover $\{U_\alpha\}$ on which
  \[
    h|_{U_\alpha}
    =\sum_{i=1}^q a_{\alpha i}g_i|_{U_\alpha},
  \]
  and choose a smooth partition of unity $\{\rho_\alpha\}$ subordinate to
  this cover. Since $\operatorname{supp}(\rho_\alpha)\subseteq U_\alpha$, the
  product $\rho_\alpha a_{\alpha i}$ extends by zero to a smooth function on
  $M$. Local finiteness then gives smooth functions
  \[
    a_i=\sum_\alpha\rho_\alpha a_{\alpha i}\in\Cinfty(M).
  \]
  Since $\sum_\alpha\rho_\alpha=1$, we obtain
  \[
    h
    =\sum_\alpha\rho_\alpha h
    =\sum_{i=1}^q a_i g_i.
  \]
\end{proof}

\begin{theorem}[Regular equations]
  \label[theorem]{thm:Regular_Equations}
  Let $g=(g_1,\ldots,g_q)\colon M\to\R^q$, and let $Z=g^{-1}(0)$. Suppose
  that
  \[
    D_pg\colon T_pM\longrightarrow\R^q
  \]
  is surjective for every $p\in Z$. Then
  \[
    \mathfrak m_Z=(g_1,\ldots,g_q).
  \]
\end{theorem}

This is \citet[Chapter~I, Lemma~2.1]{Moerdijk_Reyes_Smooth_Infinitesimal_Analysis}.

\begin{proof}
  Let $h\in\mathfrak m_Z$. We prove ideal membership locally and apply
  \Cref{lem:Local_Ideal_Membership}.

  Suppose first that $p\notin Z$. Some $g_i(p)$ is nonzero, so after shrinking
  to a neighborhood $U$ of $p$, the function $g_i|_U$ is invertible. Hence
  \[
    h|_U=(h/g_i)g_i|_U.
  \]

  Now suppose that $p\in Z$. The submersion theorem gives local coordinates
  \[
    (x_1,\ldots,x_q,y_1,\ldots,y_{d-q})
  \]
  centered at $p$ in which $g_i=x_i$ for $1\leq i\leq q$. After shrinking,
  we may take the coordinate neighborhood to be a product which is convex in
  the $x$-variables. In these coordinates, $h(0,y)=0$. The parameterized form
  of Hadamard's lemma gives
  \[
    h(x,y)
    =h(0,y)+\sum_{i=1}^q x_i
      \int_0^1\frac{\partial h}{\partial x_i}(tx,y)\,dt
    =\sum_{i=1}^q x_i h_i(x,y)
  \]
  for smooth functions $h_i$. Thus $h$ belongs locally to the ideal generated
  by the $g_i$. The local ideal-membership lemma finishes the proof.
\end{proof}

The theorem says more than the regular-value theorem. The latter identifies
$Z$ as a submanifold; the former identifies the entire ideal of functions
vanishing on $Z$.

\begin{corollary}[Function ring of a regular zero locus]
  \label[corollary]{cor:Function_Ring_Regular_Zero_Locus}
  Under the hypotheses of \Cref{thm:Regular_Equations}, restriction induces an
  isomorphism
  \[
    \Cinfty(M)/(g_1,\ldots,g_q)\cong\Cinfty(Z).
  \]
\end{corollary}

\begin{proof}
  The zero locus $Z$ is a closed embedded submanifold of $M$. Every smooth
  function on $Z$ extends to $M$: extend it along a tubular-neighborhood
  retraction, multiply by a cutoff supported in that neighborhood, and extend
  by zero. Hence the restriction map $\Cinfty(M)\to\Cinfty(Z)$ is surjective.
  Its kernel is $\mathfrak m_Z$, which is $(g_1,\ldots,g_q)$ by
  \Cref{thm:Regular_Equations}.
\end{proof}

The regularity hypothesis is essential. For $g(t)=t^2$, the zero locus is the
origin, but
\[
  (t^2)\subsetneq(t)=\mathfrak m_{\{0\}}.
\]
Consequently,
\[
  \Cinfty(\R)/(t^2)\not\cong\Cinfty(\{0\}).
\]
This is exactly the defect seen in the tangent
line--parabola intersection.

\section{Manifold function rings are finitely presented}
\label[section]{sec:Finite_Presentation_And_Products}

Our first application is a finiteness theorem: every manifold function ring
can be described by finitely many generators and smooth relations. Here
finite presentation is understood in the sense of
\Cref{def:Cinfty_Finiteness}.

\begin{example}[The circle]
  \label[example]{ex:Circle_Finite_Presentation}
  The circle has the finite presentation
  \[
    \Cinfty(S^1)\cong\Cinfty(\R^2)/(x^2+y^2-1).
  \]
  The gradient $(2x,2y)$ is nonzero on the circle, so this follows from
  the regular-equation theorem.
\end{example}

For a general manifold, a short global list of regular equations need not
be available. We will instead use a tubular neighborhood and a retraction.

The first step is to present the function ring of an arbitrary open subset
of Euclidean space. One extra generator and one relation suffice,
regardless of the shape of the open subset.

\begin{proposition}[One-relation presentation of an open subset]
  \label[proposition]{prop:Open_Subset_One_Relation}
  Let $U\subseteq\R^n$ be open. Choose a smooth function
  $a\colon\R^n\to[0,\infty)$ such that $U=\{a\neq0\}$. Then
  \[
    \Cinfty(U)
    \cong
    \Cinfty(\R^{n+1})/(ya(x)-1).
  \]
\end{proposition}

This is \citet[Chapter~I, Corollary~2.2]{Moerdijk_Reyes_Smooth_Infinitesimal_Analysis}.

\begin{proof}
  The map
  \[
    U\longrightarrow\R^{n+1},
    \qquad
    x\longmapsto\bigl(x,1/a(x)\bigr)
  \]
  is a diffeomorphism onto the hypersurface
  \[
    H_a=\{(x,y)\mid ya(x)-1=0\}.
  \]
  Along $H_a$, the partial derivative of $ya(x)-1$ with respect to $y$ is
  $a(x)$, which is nonzero. Thus the defining equation is regular, and
  \Cref{cor:Function_Ring_Regular_Zero_Locus} gives
  \[
    \Cinfty(U)\cong\Cinfty(H_a)
    \cong\Cinfty(\R^{n+1})/(ya(x)-1).
  \]
\end{proof}

A properly embedded manifold is a retract of an open tubular neighborhood
in Euclidean space. On function rings, it is therefore enough to know that
finite presentation passes to retracts.

\begin{lemma}[Retracts of finitely presented rings]
  \label[lemma]{lem:Retracts_Finitely_Presented_Cinfty_Rings}
  A retract of a finitely presented $\Cinfty$-ring is finitely presented.
\end{lemma}

\begin{proof}
  Let $A$ be a retract of
  \[
    P=\Cinfty(\R^n)/(r_1,\ldots,r_k),
  \]
  with maps $s\colon A\to P$ and $p\colon P\to A$ satisfying
  $p\circ s=\operatorname{id}_A$. The composite $e=s\circ p$ is an
  idempotent endomorphism of $P$ whose image is $s(A)$. Choose smooth
  functions $e_i\in\Cinfty(\R^n)$ representing the images under $e$ of the
  coordinate generators $x_i$ of $P$. Then $A$ is the coequalizer of
  $\operatorname{id}_P$ and $e$, and hence
  \[
    A\cong
    \Cinfty(\R^n)/
    (r_1,\ldots,r_k,x_1-e_1,\ldots,x_n-e_n).
  \]
  Indeed, a homomorphism $u\colon P\to D$ factors through $p$ precisely when
  $u=u\circ e$, and it is enough to check this equality on the coordinate
  generators. The displayed quotient therefore has the universal property of
  $A$.
\end{proof}

\begin{theorem}[Finite presentation for manifolds]
  \label[theorem]{thm:Manifold_Function_Ring_Finitely_Presented}
  If $M$ is a finite-dimensional, second-countable smooth manifold without
  boundary, then $\Cinfty(M)$ is a finitely presented $\Cinfty$-ring.
\end{theorem}

This is \citet[Chapter~I, Theorem~2.3]{Moerdijk_Reyes_Smooth_Infinitesimal_Analysis}.

\begin{proof}
  Choose a proper smooth embedding $i\colon M\hookrightarrow\R^N$. Its image
  is closed. The tubular-neighborhood theorem supplies an open neighborhood
  $U\subseteq\R^N$ of $M$ and a smooth retraction $r\colon U\to M$. Since
  $r\circ i=\operatorname{id}_M$, the induced maps
  \[
    \Cinfty(M)
    \xrightarrow{r^*}
    \Cinfty(U)
    \xrightarrow{i^*}
    \Cinfty(M)
  \]
  have composite the identity. By \Cref{prop:Open_Subset_One_Relation}, the
  middle ring is finitely presented. The result follows from
  \Cref{lem:Retracts_Finitely_Presented_Cinfty_Rings}.
\end{proof}

The theorem does not say that the underlying commutative $\R$-algebra
$\Cinfty(M)$ is finitely presented. All finiteness statements take place in
the algebraic theory whose operations are arbitrary smooth functions.

\section{Products and open restrictions}
\label[section]{sec:Products_Become_Coproducts}

The same passage from open subsets to manifolds also computes products.
Once products are understood, the equation $ya=1$ identifies restriction
to the open subset $\{a\neq0\}$ with adjoining an inverse to $a$. These
two constructions will let us reduce transverse pullbacks to local
equations and then compare the calculations on overlapping charts.

For manifolds $M$ and $P$, pullback along the two projections determines a
natural comparison homomorphism
\[
  \theta_{M,P}\colon
  \Cinfty(M)\otimes_\infty\Cinfty(P)
  \longrightarrow
  \Cinfty(M\times P).
\]

\begin{lemma}[Products of open subsets]
  \label[lemma]{lem:Products_Open_Subsets}
  If $U\subseteq\R^n$ and $V\subseteq\R^m$ are open, then
  $\theta_{U,V}$ is an isomorphism.
\end{lemma}

This is \citet[Chapter~I, Lemma~2.4]{Moerdijk_Reyes_Smooth_Infinitesimal_Analysis}.

\begin{proof}
  Choose smooth functions $a$ and $b$ with
  $U=\{a\neq0\}$ and $V=\{b\neq0\}$. By
  \Cref{prop:Open_Subset_One_Relation,prop:Cinfty_Coproduct_Formulas},
  \begin{align*}
    \Cinfty(U)\otimes_\infty\Cinfty(V)
    &\cong
    \frac{\Cinfty(\R^{n+m+2})}
      {(ya(x)-1,zb(w)-1)}.
  \end{align*}
  The two displayed equations have independent differentials along their
  common zero locus, since the coefficients of $dy$ and $dz$ are $a(x)$ and
  $b(w)$ there. Their common zero locus is
  \[
    H_a\times H_b\cong U\times V.
  \]
  The regular-equation theorem therefore identifies the displayed quotient
  with $\Cinfty(U\times V)$. Under these identifications the isomorphism is
  the natural comparison map $\theta_{U,V}$.
\end{proof}

\begin{proposition}[Products of manifolds]
  \label[proposition]{prop:Products_Manifolds_Coproducts}
  For smooth manifolds $M$ and $P$, the natural map
  \[
    \theta_{M,P}\colon
    \Cinfty(M)\otimes_\infty\Cinfty(P)
    \iso
    \Cinfty(M\times P)
  \]
  is an isomorphism.
\end{proposition}

This is \citet[Chapter~I, Proposition~2.5]{Moerdijk_Reyes_Smooth_Infinitesimal_Analysis}.

\begin{proof}
  Choose proper embeddings of $M$ and $P$, open tubular neighborhoods
  $U\supseteq M$ and $V\supseteq P$, and retractions $U\to M$ and $V\to P$.
  Thus $(M,P)$ is a retract of $(U,V)$. Naturality of $\theta$ gives a
  corresponding retraction diagram in the arrow category of $\Cinfty$-rings:
  $\theta_{M,P}$ is a retract of $\theta_{U,V}$. The latter is an isomorphism
  by \Cref{lem:Products_Open_Subsets}. A retract of an isomorphism is an
  isomorphism, so the result follows.
\end{proof}

This result is the first complete answer to
\Cref{question:When_Pushout_Is_Manifold}: a product is a pullback over a point,
and its function ring is always the corresponding coproduct.

Every open subset of a manifold can be written as $\{a\neq0\}$ for some
smooth function $a$. We now extend the one-relation presentation of an
open subset of Euclidean space to this setting.

\begin{proposition}[Smooth localization]
\label[proposition]{prop:Smooth_Localization}
Let $M$ be a manifold, let $a\in\Cinfty(M)$, and put $U=\{a\neq0\}$.
Restriction exhibits $\Cinfty(U)$ as the universal $\Cinfty$-ring under
$\Cinfty(M)$ in which $a$ becomes invertible.
\end{proposition}

\begin{proof}
By the product theorem, adjoining one free smooth generator $y$ gives
\[
 \Cinfty(M)\otimes_\infty\Cinfty(\R)\cong\Cinfty(M\times\R).
\]
The quotient by $ya-1$ therefore universally makes $a$ invertible.
The function $(m,y)\mapsto ya(m)-1$ is regular along its zero set, since its
derivative in $y$ is the nonzero number $a(m)$. That zero manifold is the
graph $m\mapsto(m,1/a(m))$ over $U$. The regular-equation theorem identifies
the quotient with $\Cinfty(U)$, and its structure map with restriction.
\end{proof}

For $M=\R^n$, this is
\citet[Chapter~I, Proposition~1.6 and Corollary~2.2]
{Moerdijk_Reyes_Smooth_Infinitesimal_Analysis}.
The argument extends it to manifolds using the product calculation in
\Cref{prop:Products_Manifolds_Coproducts}.

\section{Transverse pullbacks}
\label[section]{sec:Fiber_Products_And_The_Boundary_Of_Transversality}

We now return to \Cref{question:When_Pushout_Is_Manifold}. The essential
case is the inverse image of an embedded submanifold: transversality
ensures that its local defining equations remain independent after
pullback. General transverse fiber products reduce to this case by
using the diagonal.

\begin{theorem}[Transverse inverse images become pushouts]
  \label[theorem]{thm:Transverse_Inverse_Images_Pushouts}
  Let $f\colon M\to N$ be transverse to an embedded submanifold
  $i\colon Z\hookrightarrow N$, and put $P=f^{-1}(Z)$. Then
  \[
    \begin{tikzcd}
      \Cinfty(N) \rar["f^*"] \dar["i^*"'] \drar[pushout]
        & \Cinfty(M) \dar \\
      \Cinfty(Z) \rar
        & \Cinfty(P)
    \end{tikzcd}
  \]
  is a pushout of $\Cinfty$-rings.
\end{theorem}

This is \citet[Chapter~I, Proposition~2.6]{Moerdijk_Reyes_Smooth_Infinitesimal_Analysis}.

The local calculation contains the geometric content. Suppose
$U\subseteq N$ is an open set on which $Z$ is defined by functions
$g_1,\ldots,g_q$ with independent differentials along $U\cap Z$. Set
$V=f^{-1}(U)$. Transversality implies that
\[
  g_1\circ f,\ldots,g_q\circ f
\]
have independent differentials along $V\cap P$. Hence
\Cref{cor:Function_Ring_Regular_Zero_Locus} gives
\[
  \Cinfty(U\cap Z)
  \cong\Cinfty(U)/(g_1,\ldots,g_q)
\]
and
\[
  \Cinfty(V\cap P)
  \cong
  \Cinfty(V)/(g_1\circ f,\ldots,g_q\circ f).
\]
The local function-ring square is therefore isomorphic to
\[
  \begin{tikzcd}
    \Cinfty(U) \rar["f^*"] \dar \drar[pushout]
      & \Cinfty(V) \dar \\
    \Cinfty(U)/(g_1,\ldots,g_q) \rar
      & \Cinfty(V)/(g_1\circ f,\ldots,g_q\circ f).
  \end{tikzcd}
\]
It is a pushout by \Cref{prop:Cinfty_Coproduct_Formulas}: the lower-right ring
is obtained by imposing on $\Cinfty(V)$ the images of the same equations.

This calculation also identifies the exact use of transversality. It ensures
that the pulled-back equations remain regular, so that their algebraic
quotient is the function ring of the geometric inverse image.

The local calculation proves the theorem on coordinate neighbourhoods.
Globalization uses partitions of unity and descent for functions modulo
finitely many equations. We give that argument in
\Cref{sec:Further_Transverse_Proof}; first we explain the consequence
for arbitrary transverse fibre products.

We can now answer \Cref{question:When_Pushout_Is_Manifold} for general fiber
products. Let $f\colon M\to N$ and $g\colon P\to N$ be transverse. Then
\[
  f\times g\colon M\times P\longrightarrow N\times N
\]
is transverse to the diagonal $\Delta_N$, and
\[
  (f\times g)^{-1}(\Delta_N)=M\times_NP.
\]

\begin{theorem}[Transverse pullbacks become pushouts]
  \label[theorem]{thm:Transverse_Pullbacks_Pushouts}
  For transverse smooth maps $f\colon M\to N$ and $g\colon P\to N$, the two
  squares
  \[
    \begin{tikzcd}
      M\times_NP \rar \dar \drar[pullback]
        & P \dar["g"] \\
      M \rar["f"'] & N
    \end{tikzcd}
    \qquad
    \begin{tikzcd}
      \Cinfty(N) \rar["g^*"] \dar["f^*"'] \drar[pushout]
        & \Cinfty(P) \dar \\
      \Cinfty(M) \rar
        & \Cinfty(M\times_NP)
    \end{tikzcd}
  \]
  are respectively a pullback of manifolds and a pushout of $\Cinfty$-rings.
\end{theorem}

\begin{proof}
  Apply \Cref{thm:Transverse_Inverse_Images_Pushouts} to
  $f\times g$ and $\Delta_N\subseteq N\times N$. Then use
  \Cref{prop:Products_Manifolds_Coproducts} to identify the function rings of
  $M\times P$ and $N\times N$ with the corresponding coproducts. The resulting
  universal property says precisely that a map out of
  $\Cinfty(M\times_NP)$ is a pair of maps out of $\Cinfty(M)$ and
  $\Cinfty(P)$ whose restrictions along $f^*$ and $g^*$ agree. This is the
  pushout property of the square displayed above.
\end{proof}

Combining \Cref{cor:Recovering_Manifolds,thm:Manifold_Function_Ring_Finitely_Presented,thm:Transverse_Pullbacks_Pushouts}
gives the classical structural theorem.

\begin{theorem}[Manifolds inside smooth algebra]
  \label[theorem]{thm:Manifolds_Inside_Smooth_Algebra}
  The contravariant functor
  \[
    \Cinfty(-)\colon
    (\text{smooth manifolds})\catop
    \longrightarrow
    \Cinfty\text{-}\mathrm{Ring}
  \]
  is fully faithful, takes values in finitely presented $\Cinfty$-rings, and
  sends transverse pullback squares to pushout squares.
\end{theorem}

This is \citet[Chapter~I, Theorem~2.8]{Moerdijk_Reyes_Smooth_Infinitesimal_Analysis}.

The tangency calculation in \Cref{sec:Coproducts_Impose_Equations} shows
why the transversality hypothesis cannot simply be dropped. The geometric
intersection is a point, with function ring $\R$, while the pushout is
$\Cinfty(\R)/(x^2)\cong\R[\varepsilon]/(\varepsilon^2)$.
The comparison homomorphism kills $\varepsilon$: it forgets the
infinitesimal information retained by the equations.

This pushout shows that ordinary smooth algebra already contains
intersections which are not manifolds. It does not retain all intersection
data, however: the ordinary pushout for the self-intersection
$*\to\R\leftarrow *$ is just $\R$. The excess direction seen in
\Cref{ex:Excess_Self_Intersection}
requires a derived construction. The final exercise asks us to compare
these two descriptions.
The transverse-pullback theorem proved here is the compatibility
condition for both the ordinary and the derived universal constructions
in the next chapter.

\section{Further proofs: globalizing the pushout}
\label[section]{sec:Further_Transverse_Proof}

We supply the globalization used in
\Cref{thm:Transverse_Inverse_Images_Pushouts}. The local regular-equation
calculation is already established. The remaining task is to glue maps
into a test ring, which may contain infinitesimal elements invisible
on its real points.

\subsection{Test rings and descent}
\label[subsection]{sec:Finite_Presentation_And_Test_Rings}

To verify a pushout, we must test maps into arbitrary smooth rings. The
following criterion will let us restrict to finitely presented test rings,
where we can glue functions modulo a finite list of equations.

\begin{proposition}[Filtered-colimit criterion]
  \label[proposition]{prop:Filtered_Colimit_Criterion}
  A $\Cinfty$-ring $A$ is finitely presented if and only if
  \[
    \Hom_{\Cinfty\text{-}\mathrm{Ring}}(A,-)
  \]
  preserves filtered colimits.
\end{proposition}

\begin{proof}
  Suppose first that
  $A=\Cinfty(\R^n)/(r_1,\ldots,r_k)$. For every $\Cinfty$-ring $B$,
  \[
    \Hom(A,B)
    \cong
    \bigl\{b\in B^n\mid
      \Phi_{r_1}(b)=\cdots=\Phi_{r_k}(b)=0\bigr\}.
  \]
  Filtered colimits of $\Cinfty$-rings are computed on underlying sets, and
  filtered colimits of sets commute with finite limits. The displayed functor
  therefore preserves filtered colimits.

  Conversely, every $\Cinfty$-ring is a filtered colimit of finitely presented
  ones mapping to it. Concretely, finite sets of elements give maps from
  finitely generated free rings, finite sets of relations give finitely
  presented quotients, two such stages admit a common refinement by combining
  their generators and relations, and parallel maps become equal after adding
  finitely many relations on a finite set of generators. Apply this canonical
  filtered presentation to $A$. If $\Hom(A,-)$ preserves its colimit, then
  $\operatorname{id}_A$ factors through one finitely presented stage
  $P\to A$. Thus $A$ is a retract of $P$, and
  \Cref{lem:Retracts_Finitely_Presented_Cinfty_Rings} completes the proof.
\end{proof}

For a finitely presented test ring, the following descent statement
lets us glue the local maps constructed by imposing regular equations.

\begin{lemma}[Cech descent]
  \label[lemma]{lem:Cech_Descent_Finitely_Presented_Cinfty_Ring}
  Let
  \[
    B=\Cinfty(\R^k)/(h_1,\ldots,h_p),
    \qquad
    X=Z(h_1,\ldots,h_p),
  \]
  and let $\{W_\alpha\}$ be a cover of $X$ by open subsets of $\R^k$. For a
  nonempty finite set $S$ of indices, put
  \[
    W_S=\bigcap_{\alpha\in S}W_\alpha,
    \qquad
    B_S=
    \Cinfty(W_S)/(h_1|_{W_S},\ldots,h_p|_{W_S}).
  \]
  Then restriction induces an isomorphism
  \[
    B\iso\varprojlim_{S}B_S,
  \]
  where the limit is taken over the Cech diagram of nonempty finite
  intersections.
\end{lemma}

This is \citet[Chapter~I, Lemma~2.7]{Moerdijk_Reyes_Smooth_Infinitesimal_Analysis}.

\begin{proof}
  It is enough to prove the usual gluing statement. Suppose that we have
  $g_\alpha\in\Cinfty(W_\alpha)$ such that on every pairwise intersection
  \[
    g_\alpha-g_\beta
    \in(h_1,\ldots,h_p).
  \]
  We must construct a unique class $[g]\in B$ whose restriction is
  $[g_\alpha]$ for every $\alpha$.

  Add $\R^k\setminus X$ to the cover and choose a locally finite refinement
  $\{V_\lambda\}$ of the resulting cover of $\R^k$. If $V_\lambda$ lies in
  some $W_\alpha$, let $k_\lambda=g_\alpha|_{V_\lambda}$, choosing one such
  $\alpha$. If $V_\lambda$ lies in $\R^k\setminus X$, let $k_\lambda=0$.
  The choices are compatible modulo $(h_1,\ldots,h_p)$. Indeed, this follows
  from the original compatibility on overlaps of the $W_\alpha$, while on an
  open set disjoint from $X$ the function $h_1^2+\cdots+h_p^2$ is invertible,
  so $(h_1,\ldots,h_p)$ is the unit ideal.

  Let $\{\rho_\lambda\}$ be a partition of unity subordinate to the refined
  cover and define
  \[
    g=\sum_\lambda\rho_\lambda k_\lambda,
  \]
  extending each summand by zero. This is a globally defined smooth function.
  Fix an original index $\beta$. Near any point of $W_\beta$, only finitely
  many $\rho_\lambda$ are nonzero, so
  \[
    g-g_\beta
    =\sum_\lambda\rho_\lambda(k_\lambda-g_\beta)
  \]
  belongs locally to $(h_1,\ldots,h_p)$. By
  \Cref{lem:Local_Ideal_Membership}, it belongs to that ideal on
  $W_\beta$. Thus $[g]$ restricts to $[g_\beta]$.

  If $g'$ is another choice, then $g-g'$ belongs locally to
  $(h_1,\ldots,h_p)$ on a cover of $\R^k$. Another application of
  \Cref{lem:Local_Ideal_Membership} shows that $g-g'$ belongs to the global
  ideal. Hence $[g]=[g']$ in $B$. The construction respects every smooth
  operation, since those operations are computed componentwise in the limit.
\end{proof}

We also record how a homomorphism into a finitely presented test ring restricts
to an open cover. This is the bridge between the local geometric calculation
and \Cref{lem:Cech_Descent_Finitely_Presented_Cinfty_Ring}.

\begin{lemma}[Restricting a test-ring map]
  \label[lemma]{lem:Restricting_Test_Ring_Map}
  Let $\eta\colon\Cinfty(M)\to B$ be a homomorphism, where
  $B=\Cinfty(\R^k)/(h_1,\ldots,h_p)$, and let $\{V_\alpha\}$ be an open cover
  of $M$. There is an open cover $\{W_\alpha\}$ of
  $X=Z(h_1,\ldots,h_p)$ such that $\eta$ induces compatible homomorphisms
  \[
    \eta_S\colon
    \Cinfty(V_S)\longrightarrow B_S
  \]
  on every nonempty finite intersection.
\end{lemma}

\begin{proof}
  Choose $a_\alpha\in\Cinfty(M)$ with
  $V_\alpha=\{a_\alpha\neq0\}$. Represent $\eta(a_\alpha)$ by a smooth
  function $\widetilde a_\alpha$ on $\R^k$, and set
  \[
    W_\alpha
    =\{\widetilde a_\alpha\neq0\}.
  \]
  These open subsets of $\R^k$ cover $X$. Indeed, evaluation at $x\in X$ gives a real point
  $B\to\R$, and its composite with $\eta$ is evaluation at a point
  $m_x\in M$. Since the $V_\alpha$ cover $M$, some $a_\alpha(m_x)$ is
  nonzero, and this value equals $\widetilde a_\alpha(x)$.

  In $B_S$, each image of $\eta(a_\alpha)$ for $\alpha\in S$ is invertible.
  The universal property of adjoining smooth inverses in
  \Cref{prop:Smooth_Localization} extends
  $\eta$ uniquely to $\Cinfty(V_S)\to B_S$. Uniqueness makes these extensions
  compatible under restriction.
\end{proof}

\subsection{Globalizing the local calculation}
\label[subsection]{sec:Transverse_Inverse_Image_Theorem}

\begin{proof}[Proof of \Cref{thm:Transverse_Inverse_Images_Pushouts}]
  We first reduce to a closed embedding. An embedded submanifold $Z$
  is locally closed in $N$, so it is closed in some open neighbourhood
  $N'\subseteq N$. The inclusions $N'\to N$ and $f^{-1}(N')\to M$
  give smooth localizations. Their function-ring square is a pushout
  by the universal property of localization. Pasting with the square
  over $N'$ therefore reduces the assertion to the case where $Z$
  is closed in $N$.


  For a $\Cinfty$-ring $B$, consider the natural map
  \begin{align*}
    \Hom(\Cinfty(P),B)
    \longrightarrow
    \Hom(\Cinfty(M),B)
    \mathbin{\times_{\Hom(\Cinfty(N),B)}}
    \Hom(\Cinfty(Z),B).
  \end{align*}
  All four manifold function rings are finitely presented by
  \Cref{thm:Manifold_Function_Ring_Finitely_Presented}. Both sides therefore
  preserve filtered colimits in $B$, by
  \Cref{prop:Filtered_Colimit_Criterion} and the fact that filtered colimits of
  sets commute with finite limits. Since every $\Cinfty$-ring is a filtered
  colimit of finitely presented ones, it is enough to prove that the displayed
  map is a bijection when
  \[
    B=\Cinfty(\R^k)/(h_1,\ldots,h_p).
  \]

  Let homomorphisms
  \[
    \varphi\colon\Cinfty(M)\to B,
    \qquad
    \psi\colon\Cinfty(Z)\to B
  \]
  agree after restriction from $\Cinfty(N)$. Choose an open cover
  $\{U_\alpha\}$ of $N$ such that $U_\alpha\cap Z$ is cut out in $U_\alpha$
  by independent functions. Put $V_\alpha=f^{-1}(U_\alpha)$. The local
  calculation in \Cref{sec:Fiber_Products_And_The_Boundary_Of_Transversality} shows that the
  function-ring square over every nonempty finite intersection $U_S$ is a
  pushout.

  Choose $b_\alpha\in\Cinfty(N)$ with
  $U_\alpha=\{b_\alpha\neq0\}$. Then $V_\alpha$ is the cozero locus of
  $b_\alpha\circ f$, while $U_\alpha\cap Z$ is the cozero locus of
  $b_\alpha|_Z$. The compatibility of $\varphi$ and $\psi$ says that these two
  elements have the same image in $B$. Consequently,
  \Cref{lem:Restricting_Test_Ring_Map} produces one open cover
  $\{W_\alpha\}$ of $X=Z(h_1,\ldots,h_p)$ on which both maps restrict. We
  obtain compatible
  local maps
  \[
    \Cinfty(V_S)\longrightarrow B_S,
    \qquad
    \Cinfty(U_S\cap Z)\longrightarrow B_S.
  \]
  The local pushout property gives a unique homomorphism
  \[
    \chi_S\colon\Cinfty(V_S\cap P)\longrightarrow B_S.
  \]
  These homomorphisms are compatible on intersections, by their uniqueness.

  For $a\in\Cinfty(P)$, apply $\chi_S$ to the restriction of $a$ to
  $V_S\cap P$. The resulting elements form a compatible family in the Cech
  diagram $\{B_S\}$. By
  \Cref{lem:Cech_Descent_Finitely_Presented_Cinfty_Ring}, they determine a
  unique element $\chi(a)\in B$. This construction commutes with every smooth
  operation, so it defines a $\Cinfty$-homomorphism
  \[
    \chi\colon\Cinfty(P)\longrightarrow B.
  \]
  Locally, $\chi$ has the required composites with the restriction maps from
  $M$ and $Z$; descent makes these equalities global. The same descent argument
  shows uniqueness. Thus the displayed map of Hom-sets is a bijection.
\end{proof}

\section{Further calculations}
\label[section]{sec:Manifolds_As_Cinfty_Summary_And_Further_Reading}

\begin{example}[Smooth coproducts versus algebraic tensor products]
The difference between a smooth coproduct and an algebraic tensor product
is visible in one function. The image of
\[
 \Cinfty(\R)\otimes_\R\Cinfty(\R)\longrightarrow\Cinfty(\R^2)
\]
consists of finite sums $\sum_jf_j(x)g_j(y)$. The function $e^{xy}$
cannot be written in this form.

Indeed, if it were a sum of $r$ such terms, choose $r+1$ distinct real
numbers $x_0,\ldots,x_r$. All functions $e^{x_i y}$ would lie in the
span of $g_1,\ldots,g_r$ and hence be linearly dependent.
But differentiating a linear relation at $y=0$ in orders $0,\ldots,r$
gives a Vandermonde system with invertible matrix $(x_i^j)$.
The functions are therefore linearly independent, a contradiction.
The smooth coproduct includes $e^{xy}$ because it permits smooth
substitution involving variables from both factors.
\end{example}

\begin{exercise}
Give a finite presentation of $\Cinfty(\R\setminus\{0\})$ and use
it to calculate its coproduct with $\Cinfty(S^1)$.
\end{exercise}

\begin{exercise}
Show that the pushout for the self-intersection $*\to\R\leftarrow*$
in ordinary smooth rings is $\R$. Compare the result with the
linearized matching complex $[0\to\R]$ of \Cref{chap:Why_Derived_Manifolds}.
\end{exercise}

The proofs follow
\cite[Chapter~I, Section~2]{Moerdijk_Reyes_Smooth_Infinitesimal_Analysis}.
The differential-topological ingredients are standard tubular
neighbourhood, partition-of-unity, and submersion results; see
\cite{lee2000smooth}.
The transverse compatibility becomes part of the universal property in
\cite{Carchedi_Steffens_Universal_Derived_Manifolds}.
